\chapter{Endocarditis}

\section*{Definition and Overview}
Endocarditis is an infection of the inner lining of the heart, including the heart valves. It usually occurs when bacteria enter the bloodstream and settle on a damaged heart valve or a congenital heart defect, forming infected growths called vegetations. These growths can damage the valves, spread infection, and break off to block blood vessels elsewhere in the body. Endocarditis is a serious, life-threatening condition that requires prolonged antibiotic treatment and, in some cases, surgery. Early recognition is essential because the condition can cause lasting heart damage.

\section*{Epidemiology}
Endocarditis is uncommon but serious, affecting a small number of people per 100,000 each year. Risk is highest in people with prosthetic heart valves, previous endocarditis, congenital heart disease, and rheumatic heart disease, as well as in people who inject drugs. It is more common with age and in men. Despite modern antibiotics, endocarditis still carries a significant death rate, especially when diagnosed late.

\section*{Aetiology and Pathophysiology}
Endocarditis develops when organisms in the bloodstream settle on the heart's lining or valves.
\begin{itemize}
    \item \textbf{Bacteria:} staphylococci, streptococci, and enterococci are the most common causes
    \item \textbf{Bacteraemia:} bacteria enter the blood through dental procedures, skin infections, injections, and catheters
    \item \textbf{Predisposing conditions:} damaged or prosthetic valves, congenital heart defects, and previous endocarditis
    \item \textbf{Injecting drug use:} bacteria introduced directly into the bloodstream, often affecting the right side of the heart
    \item \textbf{Vegetations:} clumps of organisms and clotting material form on the valves
    \item \textbf{Consequences:} valve destruction, heart failure, emboli to the brain, lungs, and limbs, and abscesses in the heart
\end{itemize}

\section*{Clinical Features}
Endocarditis can present acutely or develop slowly over weeks.
\begin{itemize}
    \item Fever, chills, and sweats
    \item Fatigue, weakness, and loss of appetite
    \item Weight loss
    \item A new or changing heart murmur
    \item Small skin spots and tiny red spots in the eyes or under the nails
    \item Painful nodules in the fingers or toes
    \item Symptoms of emboli: stroke, pain in a limb, chest pain, or blood in the urine
    \item Signs of heart failure, including breathlessness and swelling
\end{itemize}

\section*{Differential Diagnosis}
Other conditions can mimic endocarditis.
\begin{itemize}
    \item Other infections, including septicaemia, pneumonia, and urinary infection
    \item Rheumatic fever and other causes of heart valve problems
    \item Connective tissue diseases, including lupus
    \item Atrial myxoma, a rare heart tumour
    \item Fever of unknown cause, with many possibilities
\end{itemize}

\section*{When to Seek Medical Care}
Anyone with persistent fever and risk factors for endocarditis, a new heart murmur, or symptoms of emboli should seek urgent medical assessment. People with known valve disease should mention their risk when unwell.

\textbf{Call 000 (or local emergency services) for:}
\begin{itemize}
    \item Symptoms of stroke: sudden weakness, facial drooping, or difficulty speaking
    \item Sudden severe chest pain or breathlessness
    \item Collapse or severe dizziness
    \item High fever with confusion or reduced consciousness
\end{itemize}

\section*{Diagnosis and Investigations}
Diagnosis is based on clinical features, blood tests, and imaging.
\begin{itemize}
    \item \textbf{Blood cultures:} the key test, taken before antibiotics where possible, to identify the organism
    \item \textbf{Echocardiography:} transthoracic and transoesophageal echo detect vegetations and valve damage
    \item \textbf{Blood tests:} markers of inflammation, blood counts, and kidney function
    \item \textbf{Urine tests:} for blood and protein
    \item \textbf{Other imaging:} CT or MRI for emboli and complications
\end{itemize}

\section*{Management}
Endocarditis requires prolonged, intensive treatment.
\begin{itemize}
    \item \textbf{Intravenous antibiotics:} weeks of high-dose antibiotics, guided by the organism's sensitivities
    \item \textbf{Hospital care:} treatment usually begins in hospital, with monitoring
    \item \textbf{Surgery:} valve repair or replacement for severe valve damage, heart failure, or uncontrolled infection
    \item \textbf{Treatment of complications:} management of emboli, abscesses, and heart failure
    \item \textbf{Long-term follow-up:} monitoring for recurrence and valve problems
\end{itemize}

\section*{Complications}
\begin{itemize}
    \item Permanent valve damage and heart failure
    \item Emboli causing stroke, kidney damage, and limb loss
    \item Abscesses in the heart muscle
    \item Kidney inflammation and failure
    \item Septic shock
    \item Recurrence of the infection
    \item Death in a significant proportion of cases
\end{itemize}

\section*{Prognosis and Outcomes}
With early diagnosis and appropriate treatment, most people with endocarditis survive, but the condition carries a real risk of death and lasting heart damage. Outcomes are best when treatment starts early and when the organism is identified. Surgery improves survival in complicated cases. After recovery, people need long-term follow-up and may need antibiotics before future procedures.

\section*{Prevention and Self-Care}
\begin{itemize}
    \item Practise good dental hygiene and have regular dental care
    \item People at high risk may need antibiotics before certain dental procedures, as advised by their doctor
    \item Never share needles, and seek help for injecting drug use
    \item Care for skin wounds to prevent infection
    \item Report fevers and new symptoms promptly if you have valve disease
    \item Carry information about your heart condition
\end{itemize}

\section*{Sources and Further Reading}
The following reputable sources provide further information for healthcare students and patients seeking reliable, current advice about endocarditis.
\begin{itemize}
    \item Heart Foundation Australia. \textit{Infective endocarditis information.} https://www.heartfoundation.org.au/
    \item Healthdirect Australia. \textit{Endocarditis.} https://www.healthdirect.gov.au/endocarditis
    \item NHS. \textit{Endocarditis: overview and treatment.} https://www.nhs.uk/conditions/endocarditis/
    \item Mayo Clinic. \textit{Endocarditis: symptoms and causes.} https://www.mayoclinic.org/diseases-conditions/endocarditis/
    \item MedlinePlus. \textit{Endocarditis.} https://medlineplus.gov/endocarditis.html
    \item MSD Manual. \textit{Infective endocarditis.} https://www.msdmanuals.com/
\end{itemize}
